\section{Corrections, integration notes, and further research}
\label{sec:corrections}

\subsection{Proposed corrections to the pinned manuscript}

The following corrections are scoped to the pinned source, so that
a maintainer can distinguish a stale manuscript statement from a gap
in the repository as a whole. Some relevant identities are already
present in existing continuation material or its exact coefficient
code. The proposed patch updates a few local statements; this article
supplies the fuller proofs and context.

\paragraph{C1. The cubic moment has a named evaluation.}
Chapter~7, at \texttt{integral:neg:cubic}, calls the identification of a
named depth-two atom an open problem. Bailey--Borwein--Borwein
\cite[Theorem~6]{BBB} already provide a Tornheim--Witten derivative
evaluation. Replace that description by the formula in
Section~\ref{sec:cubic}, or a reference to it, and state the narrower
remaining problem: reduction of the prescribed derivative combination
to a specifically defined smaller basket. The cubic-moment discovery
paragraph in Chapter~10 should make the same distinction.
A failed numerical search in a smaller basket establishes no
nonmembership statement.

\paragraph{C2. Two weight-five coordinates suffice.}
The paragraph before \texttt{gauss:eq:S4-closed} in Chapter~4 infers
only an at-most-three spanning statement from its sporadic relation.
Proposition~\ref{prop:two-g} supplies two independent exact shuffle rows,
so at most two \(g\)-coordinates remain over the displayed single
products. The sporadic relation is a proved consequence. Its proof
status and its rank consequence should both be updated, while the
mixed \(S_4\) equation itself remains conjectural.

\paragraph{C3. The imaginary weight-six mixed survivors disappear.}
The statements around \texttt{mixed:sec:eisdim} and
\texttt{mixed:thm:parallel} propose imaginary weight-six mixed
directions at both levels. Equations~\eqref{eq:mixed-weight-six}
and~\eqref{eq:eisenstein-weight-six} give explicit reductions.
The all-weight inverse-argument coefficient function in the existing
report-10 code
\texttt{10-exact-structure-verify\_mixed\_and\_euler.py}
already contains the relevant reduction mechanism.
These manuscript survivor counts should therefore be revised.
The real weight-five question is unaffected by this argument and
should be treated separately.

\paragraph{C4. A coefficient label in the sixth-point proof.}
In Chapter~7's proof of \texttt{stieltjes:\allowbreak thm:\allowbreak G1half}, the text
describes
\(\log^2 6-\log^2 3-\log^2 2\) as a coefficient of \(\zeta''(0)\).
It is the coefficient of \(\zeta(0)\).
Indeed, for \(P(s)=6^s-3^s-2^s+1\),
\[
 P(0)=P'(0)=0,\qquad
 (P\zeta)''(0)=P''(0)\zeta(0)
       =(\log^2 6-\log^2 3-\log^2 2)\zeta(0).
\]
The displayed final identity in that proof is unchanged.

\paragraph{C5. Scope of dimension and independence language.}
All matrix-rank statements should specify the formal space, background
basket, convergence or regularization rules, and exact row family.
Theorems~\ref{thm:level-four} and~\ref{thm:level-three} provide such
a specification for one imaginary quotient. They do not validate
claims about the dimensions of all actual periods. Likewise,
non-reduction in a bounded-height search should be recorded as a
search result with the basket and bounds, not as a theorem of
arithmetic independence. The existing manuscript often gives this
warning; the surviving stronger wording should be made consistent.

\paragraph{C6. An auxiliary sign in a cited preprint.}
The June 2012 author-hosted version of \cite{BBB}, equations~(63)
and~(65), prints a minus sign before a manifestly positive quadratic
sum. Section~\ref{sec:cubic} gives a direct sign check and the corrected
plus sign. This observation is restricted to the inspected preprint.
It should not be generalized to an uninspected journal version.

\subsection{Integration and claim status}

A practical integration sequence is to update the short correction
paragraphs first, then add the quantitative moment and Herglotz
sections as separate research continuations, and finally replace
finite experimental dimension summaries by explicitly scoped theorems
where applicable. The selected \(S_6\) equation belongs in the
conjecture ledger with its frozen integer vector and exact residual
certificate.

The machine-readable \texttt{claim\_ledger.json} distinguishes
analytic theorems, exact formal-algebra theorems, classical
rederivations, source corrections, and unresolved conjectures.
The package excludes exploratory rejected candidates and does not
turn a small residual into an accepted identity. All paths needed
to rebuild the article are relative, and no repository-specific
absolute path is required by the LaTeX source.

\subsection{A proposed research programme}

The following questions have concrete starting points and distinct
success criteria. Their order reflects likely usefulness rather than
a claim that any is easy.

\begin{enumerate}[leftmargin=*,label=\textbf{R\arabic*.},itemsep=0.65em]

\item \textbf{Prove the two even-index harmonic reductions.}
Find a functional or iterated-integral relation proving
Conjecture~\ref{conj:S4} and then Conjecture~\ref{conj:S6}.
The separating functional shows that re-solving the same
single-product matrix cannot suffice. Candidate routes include
relations lifted through higher depth and the extra symmetries of
the level-four punctured sphere. A useful certificate would give
the new relation explicitly and show that its image under
\(\ell\) is nonzero before eliminating it with the other rows.
Zhao~\cite{Zhao} supplies relevant broader context, but the current
obstruction does not prove that any one of those routes is necessary.

\item \textbf{Determine an even-index family with a justified basket.}
For \(S_{2m}\), decide which \(g\)-coordinates and single products
are actually sufficient at each weight. The \(S_4\) and \(S_6\)
equations suggest structured reductions, but two examples do not
determine a general coefficient law. Choose the basket from exact
relations, freeze it, record a height bound before each search, and
validate any candidate with the rational Euler method.
The desired outcome is a uniform formula or a precise counterexample
to a proposed family, rather than an unqualified claim that odd
weights require new constants.

\item \textbf{Extend the formal quotient to other levels.}
Compute the analogue of \(V_{L,w}\) first for \(L=5\) and \(6\),
with exactly the same row conventions. Determine whether its
polynomial description comes from a finite group, a module over an
invariant ring, or a more complicated relation. An all-weight Hilbert
series should be accompanied by explicit spanning and separating
families, not inferred from a short table.

\item \textbf{Compare restricted and larger relation presentations.}
Add depth-two distribution, weight-one relations, and lifted
higher-depth rows in a documented order. At each step compute the
kernel induced on the polynomial parameters of
Theorems~\ref{thm:level-four} and~\ref{thm:level-three}.
This would show exactly which new rows remove which obstructions.
A later comparison with motivic bounds must keep formal dimension,
period dimension, and conditional period conjectures distinct.

\item \textbf{Uniform moment truncation beyond the least-term window.}
Theorem~\ref{thm:optimal-moment} treats bounded
\(\eta=N\Lambda-n\). Develop a uniform expansion when \(N/n\)
varies over a positive interval, and then when it tends to zero or
infinity. The inverse remainder kernel and gamma concentration
suggest the relevant starting formula. The challenge is to retain
uniform error estimates through the changing concentration region.
A second target is a fully explicit finite-\(n,N\) remainder bound.

\item \textbf{Global analytic structure of the inverse gamma branch.}
Locate the next singularities reached by analytic continuation of
the particular branch \(g\), including which critical values
actually lie on that branch. Determine their contributions after
the leading square-root expansion is removed. The local slit-disk
argument used here deliberately avoids an unsupported assertion
about a global inverse. Understanding the further singularities
could explain exponentially smaller corrections to the late
coefficients and constrain possible global sign bounds.

\item \textbf{Higher Herglotz corrections and transition in the argument.}
Compute several more half-integer terms in the oscillatory expansion
and give explicit error constants. More substantially, analyze
simultaneous limits such as \(x/r\to\xi>0\) or \(x\to0\) with \(r\).
The fixed-\(x\), compact-uniform theorem does not automatically cover
these regimes: the optimal Mellin contour and the contributing
arithmetic sectors change. A useful result would identify the
transition from one dominant sector to a collective contribution.

\item \textbf{Quantitative Herglotz sign locations.}
Theorem~\ref{thm:herglotz-oscillation} gives infinitely many signs,
but does not identify every sign for every large integer \(r\).
Use higher corrections and explicit bounds to locate the changes
near the zeros of the leading cosine. Because the phase increment
tends to zero, an additive error estimate is essential near these
zeros. The target is an eventually certified sign algorithm and an
asymptotic count of transitions.

\item \textbf{Arithmetic twists of the high-derivative theorem.}
Replace the divisor coefficients by character-twisted coefficients
or consider higher Herglotz functions. The functional-equation
framework of \cite{RZ,DSS} is a natural starting point.
Determine how parity, conductor, and vanishing of the first
arithmetic coefficient alter the exponential rate and oscillatory
phase. The sum of all sectors must again be bounded uniformly;
a single special-function asymptotic is not enough.

\item \textbf{Minimal derivative coordinates for the cubic moment.}
The cubic formula singles out \(2T_{13}-T_{12}\), together with
\(T_3=2Z'(2)\). Determine the exact relations among Tornheim
derivatives at \((1,1,1)\) obtained from symmetry, integration by
parts, and differentiated functional equations. State a candidate
basis before asking whether this combination reduces to ordinary
zeta derivatives. Failure of such a reduction search should not be
confused with the already solved problem of finding a named
depth-two representation.

\item \textbf{Certified evaluation of the asymptotic quantities.}
Combine the explicit tail estimates with interval implementations
of gamma, polygamma, Hurwitz zeta, and the remaining finite
integrals. For moments, the cutoff identity
\eqref{eq:cutoff} provides a convergent route; for Herglotz remainders,
\eqref{eq:herglotz-simple-tail} provides a simple arithmetic cutoff.
The rational \(S_6\) certificate is a model of the distinction between
a proved truncation bound and an actual certified implementation.

\item \textbf{Repository-level mathematical provenance.}
Attach to every accepted identity its exact normalization,
definition of the basis, proof source, executable check, and status.
For dimension claims also store the row generator and a basis or
separating witness. For numerical conjectures store the frozen
coefficients separately from discovery logs, with an independent
holdout computation. This would prevent a correction already
proved in a continuation from remaining an open claim in the
consolidated manuscript.

\end{enumerate}

These questions connect the analytic and algebraic sides of the
programme: good coordinates make strong conjectures visible, while
controlled remainders and explicit relation spaces show exactly
what has, and has not, been proved.

